\chapter{Parallelism and Concurrency}
\begin{quotation}\textit{
With many simple computer systems it is natural to expect that the
computer will simply follow its instructions and that if the same
input is provided on several occasions exactly the same behaviour will be
observed. Once parallelism is introduced this tidy situation falls, since
the precise timing for each of multiple concurrent parts of the activity
can not be predicted. When you have multiple tasks to perform, they are all
independent and the order in which thet complete is utterly unimportant all
is still well, but that very special situation is far from common. \\
In this chapter three styles of concurrency are considered:
\begin{enumerate}
\item Multiple-core CPUs with work spread across them;
\item ``Single Instruction Multiple Data'' parallelism as often
implemented using graphics cards;
\item Work distributed across a network of interconnected systems.
\end{enumerate}
Here it is firstly obviously essential that uncertain timing as between
parts of thw workload must never lead to the whole task getting out of step
and perhaps locking up. But then maintaining that level of synchronization
can easily lead to severe overheads, and the standard joke that the
computation being run has been paralysed as much as parallelised.
}\end{quotation}
\section{Multiple Cores}
\subsection{Software}
% Atomic variables, semaphores, condition variables, higher level constructs
\subsection{Hardware}
% memory models and acquire/release semantics to worry about, especially
% in lock-free code.
\section{SIMD}
% The sequential core of what you do may dominate costs!
\section{Distributed Computing}
% Communication costs and balancing are not easy to optimize!


